\chapter{Meiosis and Life Cycles}\label{ch:g11:meiosis-life-cycles}

\omterm{def:g11:meiosis-life-cycles:repro}{Sexual reproduction} mixes \omterm{prop:g11:dna-replication:genome}{genomes}. That mixing depends on a special division
--- \emph{\omterm{def:g11:meiosis-life-cycles:meiosis}{meiosis}} --- that halves the \omterm{def:g11:dna-replication:chromosome}{chromosome} number and reshuffles alleles
through independent assortment and \omterm{prop:g11:meiosis-life-cycles:prophaseI}{crossing over}. This chapter contrasts \omterm{def:g11:meiosis-life-cycles:meiosis}{meiosis}
with \omterm{def:g11:cell-cycle-mitosis:mitosis}{mitosis}, then places both in simple life cycles.

\section{Why meiosis?}

\begin{definition}[Haploid and diploid]\label{def:g11:meiosis-life-cycles:ploidy}
A \emph{haploid}\index{haploid} \omterm{def:g10:the-cell:cell}{cell} has one set of \omterm{def:g11:dna-replication:chromosome}{chromosomes} ($n$). A
\emph{diploid}\index{diploid} \omterm{def:g10:the-cell:cell}{cell} has two homologous sets ($2n$). Human gametes
are haploid ($n = 23$); a zygote formed by fertilization is diploid ($2n = 46$).
\end{definition}

\begin{proposition}[Fertilization without meiosis would double ploidy]\label{prop:g11:meiosis-life-cycles:need}
If two \omterm{def:g11:meiosis-life-cycles:ploidy}{diploid} \omterm{def:g10:the-cell:cell}{cells} fused each generation, \omterm{def:g11:dna-replication:chromosome}{chromosome} number would double every
generation. \omterm{def:g11:meiosis-life-cycles:meiosis}{Meiosis} produces \omterm{def:g11:meiosis-life-cycles:ploidy}{haploid} gametes so that fertilization restores $2n$
rather than creating $4n$.
\end{proposition}
\admitted

\begin{definition}[Meiosis]\label{def:g11:meiosis-life-cycles:meiosis}
\emph{Meiosis}\index{meiosis} is two consecutive nuclear divisions
(meiosis~I and meiosis~II) that convert one \omterm{def:g11:meiosis-life-cycles:ploidy}{diploid} \omterm{def:g10:the-cell:cell}{cell} into four \omterm{def:g11:meiosis-life-cycles:ploidy}{haploid}
\omterm{def:g10:the-cell:cell}{cells}, after a single round of \omterm{def:g11:dna-replication:semi}{DNA replication}.
\end{definition}

\section{Meiosis I: separating homologues}

\begin{proposition}[Prophase I and crossing over]\label{prop:g11:meiosis-life-cycles:prophaseI}
In \emph{prophase I}\index{prophase I}, homologous \omterm{def:g11:dna-replication:chromosome}{chromosomes} pair
(\emph{synapsis}\index{synapsis}) as a \emph{bivalent}\index{bivalent} (tetrad of
four chromatids). \emph{Crossing over}\index{crossing over} exchanges
corresponding segments between non-sister chromatids, producing recombinant
chromatids. Visible X-shaped connections are \emph{chiasmata}\index{chiasma}.
\end{proposition}
\admitted

\begin{definition}[Meiosis I overview]\label{def:g11:meiosis-life-cycles:MI}
In \emph{meiosis I}\index{meiosis I} (reductional division):
\begin{itemize}
  \item metaphase I: \omterm{prop:g11:meiosis-life-cycles:prophaseI}{bivalents} align at the equator;
  \item anaphase I: homologous \omterm{def:g11:dna-replication:chromosome}{chromosomes} separate to opposite poles (\omterm{def:g11:cell-cycle-mitosis:chromatid}{sister
        chromatids} remain together);
  \item telophase I / \omterm{def:g11:cell-cycle-mitosis:cytokinesis}{cytokinesis}: two \omterm{def:g11:meiosis-life-cycles:ploidy}{haploid} \omterm{def:g10:the-cell:cell}{cells} form, each \omterm{def:g11:dna-replication:chromosome}{chromosome} still
        composed of two chromatids.
\end{itemize}
\end{definition}

\begin{omfigure}
\begin{tikzpicture}[font=\small]
  % bivalent
  \draw[omDef, line width=2.5pt] (0,0.3) -- (0,2.0);
  \draw[omDef, line width=2.5pt] (0.35,0.3) -- (0.35,2.0);
  \draw[omThm, line width=2.5pt] (0.9,0.3) -- (0.9,2.0);
  \draw[omThm, line width=2.5pt] (1.25,0.3) -- (1.25,2.0);
  \draw[omProp, thick] (0.35,1.2) -- (0.9,1.2);
  \node[below, align=center] at (0.6,-0.1) {bivalent\\(crossing over)};
  \node[above, omDef] at (0.15,2.05) {maternal};
  \node[above, omThm] at (1.1,2.05) {paternal};
  % arrow
  \draw[-Stealth] (2.0,1.1) -- (2.8,1.1);
  \node[above] at (2.4,1.2) {anaphase I};
  % products
  \draw[omDef, line width=2.5pt] (3.5,1.5) -- (3.5,2.5);
  \draw[omDef, line width=2.5pt] (3.85,1.5) -- (3.85,2.5);
  \draw[omThm, line width=2.5pt] (3.5,0) -- (3.5,1.0);
  \draw[omThm, line width=2.5pt] (3.85,0) -- (3.85,1.0);
  \node[right, align=left] at (4.1,2.0) {haploid\\(dyad)};
  \node[right, align=left] at (4.1,0.5) {haploid\\(dyad)};
\end{tikzpicture}

{\small \omterm{def:g11:meiosis-life-cycles:MI}{Meiosis I} separates homologues. \omterm{prop:g11:meiosis-life-cycles:prophaseI}{Crossing over} has already swapped
pieces between non-sister chromatids (chiasma region schematized).}
\end{omfigure}

\section{Meiosis II: separating sisters}

\begin{definition}[Meiosis II]\label{def:g11:meiosis-life-cycles:MII}
\emph{Meiosis II}\index{meiosis II} resembles \omterm{def:g11:cell-cycle-mitosis:mitosis}{mitosis}: \omterm{def:g11:cell-cycle-mitosis:chromatid}{sister chromatids}
separate at anaphase II. No \omterm{def:g11:dna-replication:semi}{DNA replication} occurs between \omterm{def:g11:meiosis-life-cycles:meiosis}{meiosis}~I and~II.
The four resulting nuclei are \omterm{def:g11:meiosis-life-cycles:ploidy}{haploid} with unreplicated \omterm{def:g11:dna-replication:chromosome}{chromosomes}.
\end{definition}

\begin{proposition}[Mitosis versus meiosis]\label{prop:g11:meiosis-life-cycles:vs}
\begin{center}\small
\begin{tabular}{lll}
 & \textbf{\omterm{def:g11:cell-cycle-mitosis:mitosis}{Mitosis}} & \textbf{\omterm{def:g11:meiosis-life-cycles:meiosis}{Meiosis}} \\
Rounds of \omterm{def:g11:dna-replication:semi}{DNA replication} & $1$ & $1$ \\
Nuclear divisions & $1$ & $2$ \\
Daughter \omterm{def:g10:the-cell:cell}{cells} & $2$ \omterm{def:g11:meiosis-life-cycles:ploidy}{diploid} & $4$ \omterm{def:g11:meiosis-life-cycles:ploidy}{haploid} \\
Pairing of homologues & no & yes (\omterm{def:g11:meiosis-life-cycles:MI}{meiosis I}) \\
\omterm{prop:g11:meiosis-life-cycles:prophaseI}{Crossing over} & rare / not routine & routine in \omterm{prop:g11:meiosis-life-cycles:prophaseI}{prophase I} \\
Role & \omterm{def:g10:scientific-biology:properties}{growth}, repair & gametes / spores \\
\end{tabular}
\end{center}
\end{proposition}
\admitted

\begin{omfigure}
\begin{tikzpicture}[font=\footnotesize]
  \node[draw=omDef, thick, align=center, fill=omDef!8] (p) at (0,2.2)
    {diploid parent\\$2n$, after S};
  \node[draw=omMeth, thick, align=center, fill=omMeth!8] (m) at (0,0.8)
    {two cells\\$n$ (2 chromatids)};
  \node[draw=omProp, thick, align=center, fill=omProp!8] (h1) at (-1.8,-0.6)
    {$n$};
  \node[draw=omProp, thick, align=center, fill=omProp!8] (h2) at (-0.6,-0.6)
    {$n$};
  \node[draw=omProp, thick, align=center, fill=omProp!8] (h3) at (0.6,-0.6)
    {$n$};
  \node[draw=omProp, thick, align=center, fill=omProp!8] (h4) at (1.8,-0.6)
    {$n$};
  \draw[-Stealth] (p) -- (m) node[midway, right] {meiosis I};
  \draw[-Stealth] (m) -- (h1);
  \draw[-Stealth] (m) -- (h2);
  \draw[-Stealth] (m) -- (h3);
  \draw[-Stealth] (m) -- (h4);
  \node at (2.8,0.1) {meiosis II};
\end{tikzpicture}

{\small One \omterm{def:g11:dna-replication:semi}{replication}, two divisions: four \omterm{def:g11:meiosis-life-cycles:ploidy}{haploid} products. In many animals
only one product becomes a large egg; the others are \omterm{def:g11:meiosis-life-cycles:gametogenesis}{polar bodies}.}
\end{omfigure}

\section{Sources of genetic variation}

\begin{proposition}[Independent assortment]\label{prop:g11:meiosis-life-cycles:assortment}
At metaphase I, each \omterm{prop:g11:meiosis-life-cycles:prophaseI}{bivalent} orients independently of the others. For $n$
\omterm{def:g11:dna-replication:chromosome}{chromosome} pairs there are $2^{n}$ possible combinations of maternal and
paternal \omterm{def:g11:dna-replication:chromosome}{chromosomes} in a gamete (ignoring \omterm{prop:g11:meiosis-life-cycles:prophaseI}{crossing over}). Humans:
$2^{23} \approx \num{8.4e6}$ combinations from assortment alone.
\end{proposition}
\admitted

\begin{proposition}[Crossing over multiplies diversity]\label{prop:g11:meiosis-life-cycles:co-var}
\omterm{prop:g11:meiosis-life-cycles:prophaseI}{Crossing over} creates chromatids that are mosaics of maternal and paternal
alleles. Together with independent assortment and random fertilization, it makes
siblings genetically unique (except identical twins).
\end{proposition}
\admitted

\begin{method}[Counting combinations (no crossing over)]\label{met:g11:meiosis-life-cycles:count}
For an \omterm{def:g10:scientific-biology:organism}{organism} with $n$ pairs of \omterm{def:g11:dna-replication:chromosome}{chromosomes}, number of distinct gamete
genotypes from independent assortment only is $2^{n}$. With random fertilization
the number of zygote combinations is $(2^{n})^{2} = 2^{2n}$.
\end{method}

\begin{example}[A small genome]\label{ex:g11:meiosis-life-cycles:small}
If $2n = 4$ ($n = 2$), assortment alone yields $2^{2} = 4$ gamete types and
$4 \times 4 = 16$ zygote combinations (again ignoring \omterm{prop:g11:meiosis-life-cycles:prophaseI}{crossing over}).
\end{example}

\section{Oogenesis and spermatogenesis}

\begin{definition}[Gametogenesis in animals]
\label{def:g11:meiosis-life-cycles:gametogenesis}
\emph{Spermatogenesis}\index{spermatogenesis} produces four \omterm{def:g11:meiosis-life-cycles:ploidy}{haploid} sperm from
each \omterm{def:g10:the-cell:cell}{cell} that completes \omterm{def:g11:meiosis-life-cycles:meiosis}{meiosis}. \emph{Oogenesis}\index{oogenesis} typically
produces one large \omterm{def:g11:meiosis-life-cycles:ploidy}{haploid} egg and small \emph{polar bodies}\index{polar body}
that discard extra \omterm{def:g11:meiosis-life-cycles:ploidy}{haploid} nuclei, conserving cytoplasm and resources in a
single ovum.
\end{definition}

\begin{omfigure}
\begin{tikzpicture}[font=\footnotesize,
  c/.style={draw, circle, minimum size=0.55cm, inner sep=0pt, align=center}]
  % spermatogenesis column
  \node[align=center] at (0,3.4) {\textbf{spermatogenesis}};
  \node[c, fill=omDef!20] (s0) at (0,2.6) {};
  \node[below=1pt] at (s0.south) {$2n$};
  \draw[-Stealth] (0,2.2) -- (0,1.85);
  \node[c, fill=omDef!20] (s1a) at (-0.55,1.4) {};
  \node[c, fill=omDef!20] (s1b) at (0.55,1.4) {};
  \draw[-Stealth] (-0.55,1.05) -- (-0.55,0.7);
  \draw[-Stealth] (0.55,1.05) -- (0.55,0.7);
  \node[c, fill=omProp!25] (s2a) at (-0.9,0.35) {};
  \node[c, fill=omProp!25] (s2b) at (-0.2,0.35) {};
  \node[c, fill=omProp!25] (s2c) at (0.2,0.35) {};
  \node[c, fill=omProp!25] (s2d) at (0.9,0.35) {};
  \node[align=center] at (0,-0.35) {four sperm\\($n$ each)};
  % oogenesis column
  \node[align=center] at (5.2,3.4) {\textbf{oogenesis}};
  \node[c, fill=omMeth!20, minimum size=0.75cm] (o0) at (5.2,2.55) {};
  \node[below=1pt] at (o0.south) {$2n$};
  \draw[-Stealth] (5.2,2.05) -- (5.2,1.7);
  \node[c, fill=omMeth!25, minimum size=0.7cm] (o1) at (4.7,1.2) {};
  \node[c, fill=omExo!25, minimum size=0.28cm] (pb1) at (5.9,1.2) {};
  \node[font=\scriptsize] at (5.9,0.75) {polar};
  \draw[-Stealth] (4.7,0.75) -- (4.7,0.4);
  \node[c, fill=omMeth!30, minimum size=0.75cm] (egg) at (4.5,-0.05) {};
  \node[c, fill=omExo!25, minimum size=0.25cm] (pb2) at (5.5,-0.05) {};
  \node[align=center] at (5.0,-0.75) {one egg + polar bodies};
\end{tikzpicture}

{\small Equal \omterm{def:g11:cell-cycle-mitosis:cytokinesis}{cytokinesis} in \omterm{def:g11:meiosis-life-cycles:gametogenesis}{spermatogenesis} versus unequal \omterm{def:g11:cell-cycle-mitosis:cytokinesis}{cytokinesis} in
\omterm{def:g11:meiosis-life-cycles:gametogenesis}{oogenesis}: same meiotic \omterm{def:g11:dna-replication:chromosome}{chromosome} logic, different investment per gamete.}
\end{omfigure}

\begin{proposition}[Same divisions, different packaging]
\label{prop:g11:meiosis-life-cycles:gameto-vs}
Both processes use \omterm{def:g11:meiosis-life-cycles:MI}{meiosis I} and II after one \omterm{def:g11:cell-cycle-mitosis:interphase}{S phase}. They differ in
\omterm{def:g11:cell-cycle-mitosis:cytokinesis}{cytokinesis}, timing (continuous sperm production after puberty in human males;
prolonged oocyte arrest in females) and the number of functional gametes per
\omterm{def:g11:meiosis-life-cycles:meiosis}{meiosis} --- not in the fundamental need to produce \omterm{def:g11:meiosis-life-cycles:ploidy}{haploid} nuclei.
\end{proposition}
\admitted

\section{Nondisjunction and aneuploidy}

\begin{definition}[Nondisjunction in meiosis]
\label{def:g11:meiosis-life-cycles:nondisjunction}
\emph{Nondisjunction}\index{nondisjunction} is failure of homologous \omterm{def:g11:dna-replication:chromosome}{chromosomes}
(\omterm{def:g11:meiosis-life-cycles:MI}{meiosis I}) or \omterm{def:g11:cell-cycle-mitosis:chromatid}{sister chromatids} (\omterm{def:g11:meiosis-life-cycles:MII}{meiosis II}) to separate properly. Gametes
receive $n+1$ or $n-1$ \omterm{def:g11:dna-replication:chromosome}{chromosomes}; fertilisation then yields
\emph{aneuploid}\index{aneuploidy} zygotes (e.g.\ trisomy or monosomy).
\end{definition}

\begin{remark}[Link to human genetics]
\label{rem:g11:meiosis-life-cycles:aneuploidy-link}
Trisomy 21 and \omterm{def:g11:dna-replication:chromosome}{sex-chromosome} aneuploidies (XXY, X0) arise this way
(\cref{ch:g11:human-genetics}). Risk of some \omterm{def:g11:meiosis-life-cycles:nondisjunction}{nondisjunction} events increases
with maternal age. \omterm{def:g11:meiosis-life-cycles:meiosis}{Meiosis} is therefore not only a source of healthy
variation --- errors in segregation have clinical consequences.
\end{remark}

\begin{omfigure}
\begin{tikzpicture}[font=\footnotesize]
  \node[align=center] at (0,2.0) {meiosis I error};
  \draw[omDef, thick] (-0.4,0.8) -- (-0.4,1.6);
  \draw[omDef, thick] (0,0.8) -- (0,1.6);
  \draw[omThm, thick] (-0.4,0.8) -- (-0.4,1.6);
  \draw[omThm, thick] (0,0.8) -- (0,1.6);
  \draw[-Stealth, omThm] (-0.2,1.2) -- (-0.2,1.7);
  \node[align=center] at (0,0.25) {both homologues\\to one pole};
  \node[align=center] at (4.2,2.0) {after fertilisation};
  \node[draw=omProp, thick, align=center, fill=omProp!10] at (4.2,0.9)
    {$2n+1$ or $2n-1$\\zygote};
\end{tikzpicture}

{\small \omterm{def:g11:meiosis-life-cycles:nondisjunction}{Nondisjunction} produces \omterm{def:g11:meiosis-life-cycles:nondisjunction}{aneuploid} gametes; fusion with a normal gamete
yields trisomic or monosomic zygotes.}
\end{omfigure}

\section{Sexual and asexual reproduction}

\begin{definition}[Asexual and sexual reproduction]\label{def:g11:meiosis-life-cycles:repro}
\emph{Asexual reproduction}\index{asexual reproduction} produces offspring from
one parent without fusion of gametes (binary fission, budding, vegetative
spread, parthenogenesis in some species). \emph{Sexual reproduction}\index{sexual
reproduction} involves \omterm{def:g11:meiosis-life-cycles:meiosis}{meiosis} and fertilization (or equivalent fusion of
\omterm{def:g11:meiosis-life-cycles:ploidy}{haploid} \omterm{def:g10:the-cell:cell}{cells}).
\end{definition}

\begin{remark}
Asexual lineages can multiply quickly when conditions are stable. \omterm{def:g11:meiosis-life-cycles:repro}{Sexual
reproduction} continually recombines alleles, which can help \omterm{def:g10:ecosystems:levels}{populations} track
changing environments and parasites --- a major theme of evolutionary \omterm{def:g10:scientific-biology:science}{biology}
(\cref{ch:g11:evolution-mechanisms}).
\end{remark}

\section{Life-cycle sketches}

\begin{definition}[Alternation of generations (outline)]\label{def:g11:meiosis-life-cycles:alt}
Many plants and algae show \emph{alternation of generations}\index{alternation of
generations}: a multicellular \omterm{def:g11:meiosis-life-cycles:ploidy}{diploid} \emph{sporophyte}\index{sporophyte}
produces \omterm{def:g11:meiosis-life-cycles:ploidy}{haploid} spores by \omterm{def:g11:meiosis-life-cycles:meiosis}{meiosis}; spores grow into a multicellular \omterm{def:g11:meiosis-life-cycles:ploidy}{haploid}
\emph{gametophyte}\index{gametophyte} that produces gametes by \omterm{def:g11:cell-cycle-mitosis:mitosis}{mitosis}; gametes
fuse to a zygote that grows into a new sporophyte.
\end{definition}

\begin{omfigure}
\begin{tikzpicture}[font=\small]
  \node[draw=omDef, thick, fill=omDef!10, minimum width=2.4cm] (sp)
    at (0,1.5) {sporophyte $2n$};
  \node[draw=omMeth, thick, fill=omMeth!10, minimum width=2.4cm] (ga)
    at (0,-0.8) {gametophyte $n$};
  \draw[-Stealth, omProp] (sp.south) -- (ga.north)
    node[midway, right] {meiosis $\to$ spores};
  \draw[-Stealth, omThm] (ga.east) -- ++(1.4,0) -- ++(0,2.3) -- (sp.east)
    node[pos=0.25, right] {mitosis $\to$ gametes};
  \node[right, align=left] at (1.6,0.7) {fertilization\\$\to$ zygote $2n$};
\end{tikzpicture}

{\small Plant-style alternation: \omterm{def:g11:meiosis-life-cycles:meiosis}{meiosis} makes spores (not gametes); the
\omterm{def:g11:meiosis-life-cycles:alt}{gametophyte} makes gametes by \omterm{def:g11:cell-cycle-mitosis:mitosis}{mitosis}. Animals skip a multicellular \omterm{def:g11:meiosis-life-cycles:ploidy}{haploid}
phase: \omterm{def:g11:meiosis-life-cycles:meiosis}{meiosis} makes gametes directly.}
\end{omfigure}

\begin{proposition}[Animal life cycle sketch]\label{prop:g11:meiosis-life-cycles:animal}
In animals the multicellular body is \omterm{def:g11:meiosis-life-cycles:ploidy}{diploid}. \omterm{def:g11:meiosis-life-cycles:meiosis}{Meiosis} in the germline produces
\omterm{def:g11:meiosis-life-cycles:ploidy}{haploid} gametes; fertilization restores diploidy; the zygote grows by \omterm{def:g11:cell-cycle-mitosis:mitosis}{mitosis}.
There is no multicellular \omterm{def:g11:meiosis-life-cycles:alt}{gametophyte}.
\end{proposition}
\admitted

\section{Exercises}

\begin{exercise}[$\star$]\label{exo:g11:meiosis-life-cycles:1}
Why must \omterm{def:g11:meiosis-life-cycles:meiosis}{meiosis} precede fertilization in a sexual life cycle that keeps a
constant \omterm{def:g11:dna-replication:chromosome}{chromosome} number?
\end{exercise}

\begin{exercise}[$\star$]\label{exo:g11:meiosis-life-cycles:2}
How many \omterm{def:g11:dna-replication:semi}{DNA replications} and how many nuclear divisions occur in \omterm{def:g11:meiosis-life-cycles:meiosis}{meiosis}?
\end{exercise}

\begin{exercise}[$\star$]\label{exo:g11:meiosis-life-cycles:3}
What separates at anaphase I? What separates at anaphase II?
\end{exercise}

\begin{exercise}[$\star\star$]\label{exo:g11:meiosis-life-cycles:4}
Define \omterm{prop:g11:meiosis-life-cycles:prophaseI}{crossing over}. Between which chromatids does it occur?
\end{exercise}

\begin{exercise}[$\star\star$]\label{exo:g11:meiosis-life-cycles:5}
A \omterm{def:g10:the-cell:cell}{cell} has $2n = 8$. How many \omterm{def:g11:dna-replication:chromosome}{chromosomes} are in each product of \omterm{def:g11:meiosis-life-cycles:meiosis}{meiosis}? How
many chromatids per \omterm{def:g11:dna-replication:chromosome}{chromosome} just after \omterm{def:g11:meiosis-life-cycles:meiosis}{meiosis}~I (before \omterm{def:g11:meiosis-life-cycles:meiosis}{meiosis}~II)?
\end{exercise}

\begin{exercise}[$\star\star$]\label{exo:g11:meiosis-life-cycles:6}
List three differences between \omterm{def:g11:cell-cycle-mitosis:mitosis}{mitosis} and \omterm{def:g11:meiosis-life-cycles:meiosis}{meiosis}.
\end{exercise}

\begin{exercise}[$\star\star$]\label{exo:g11:meiosis-life-cycles:7}
For $n = 3$, how many gamete \omterm{def:g11:dna-replication:chromosome}{chromosome} combinations arise from independent
assortment alone?
\end{exercise}

\begin{exercise}[$\star\star$]\label{exo:g11:meiosis-life-cycles:8}
In \omterm{def:g11:meiosis-life-cycles:alt}{alternation of generations}, which generation is \omterm{def:g11:meiosis-life-cycles:ploidy}{diploid}, and what process
produces spores?
\end{exercise}

\begin{exercise}[$\star\star\star$]\label{exo:g11:meiosis-life-cycles:9}
Explain how independent assortment and \omterm{prop:g11:meiosis-life-cycles:prophaseI}{crossing over} both increase genetic
variation among gametes, using different mechanisms.
\end{exercise}

\begin{exercise}[$\star\star\star$]\label{exo:g11:meiosis-life-cycles:10}
A student says ``\omterm{def:g11:meiosis-life-cycles:meiosis}{meiosis} is just two mitoses in a row.'' Correct the statement
with two precise objections.
\end{exercise}

\begin{exercise}[$\star\star$]\label{exo:g11:meiosis-life-cycles:11}
Compare \omterm{def:g11:meiosis-life-cycles:gametogenesis}{oogenesis} and \omterm{def:g11:meiosis-life-cycles:gametogenesis}{spermatogenesis}: number of functional gametes per
\omterm{def:g11:meiosis-life-cycles:meiosis}{meiosis}, relative size of products, and one timing difference in human
life history.
\end{exercise}

\begin{exercise}[$\star\star\star$]\label{exo:g11:meiosis-life-cycles:12}
Explain how \omterm{def:g11:meiosis-life-cycles:nondisjunction}{nondisjunction} at \omterm{def:g11:meiosis-life-cycles:MI}{meiosis I} can produce a trisomic zygote after
fertilisation by a normal gamete. Why is anaphase I (homologue separation)
a critical control point for \omterm{def:g11:dna-replication:chromosome}{chromosome} number?
\end{exercise}

\begin{problem}[Designing a meiosis diagram]\label{pb:g11:meiosis-life-cycles:1}
An \omterm{def:g10:scientific-biology:organism}{organism} has $2n = 4$ \omterm{def:g11:dna-replication:chromosome}{chromosomes}: one long pair and one short pair. Maternal
\omterm{def:g11:dna-replication:chromosome}{chromosomes} are marked M, paternal P.
\begin{enumerate}
  \item Draw the \omterm{def:g10:the-cell:cell}{cell} in G$_2$ before \omterm{def:g11:meiosis-life-cycles:meiosis}{meiosis} (four \omterm{def:g11:dna-replication:chromosome}{chromosomes} as dyads).
  \item Draw metaphase I for one possible orientation of the two \omterm{prop:g11:meiosis-life-cycles:prophaseI}{bivalents}.
  \item Draw the two \omterm{def:g10:the-cell:cell}{cells} after \omterm{def:g11:meiosis-life-cycles:meiosis}{meiosis}~I.
  \item Draw the four \omterm{def:g10:the-cell:cell}{cells} after \omterm{def:g11:meiosis-life-cycles:meiosis}{meiosis}~II for that orientation.
  \item How many different orientations of the two \omterm{prop:g11:meiosis-life-cycles:prophaseI}{bivalents} are possible at
        metaphase I? List the gamete types (by M/P content) from assortment
        alone.
  \item Indicate on one \omterm{prop:g11:meiosis-life-cycles:prophaseI}{bivalent} where a single crossover between non-sister
        chromatids would create recombinant chromatids.
  \item Contrast this sexual pathway with binary fission of a bacterium in terms
        of offspring genetic identity.
\end{enumerate}
\end{problem}
